\section{Method}

% ... (前面已有的 Method 子节，如 Overall Framework、PegaAgent 概述等) ...

\subsection{Long-chain Implicit Decision Mechanism}
\label{subsec:lidm}

\subsubsection{Motivation}

Long-horizon GUI agent tasks often require cross-app reasoning where critical
decision-making information is \textit{implicit}---not stated in the task goal,
but scattered across files, screens, or app states that the agent must
discover and read on its own. We define such problems as
\textbf{Long-chain Implicit Decision Problems} (LIDPs): tasks where the agent
must (1) navigate multiple apps to gather implicit information, (2) reason over
this information using rules that are also not part of the goal, and (3) make a
correct downstream decision based on the gathered facts.

Consider the \textit{MedicalAppointmentWorkflow} task: the agent must read a
symptom note, read a clinic booking policy (both hidden from the goal), count
the symptoms, and decide whether to book a ``Specialist Appointment''
($>3$ symptoms) or a ``Routine Checkup'' ($\leq 3$ symptoms). In our
experiments, baseline agents fail at Step 18 (filling the appointment title)
because the symptom count extracted at Step 10 is lost by the time the
decision is made. The root cause is that the agent's memory consists solely
of natural-language history summaries, which discard structured facts during
compression.

\subsubsection{Design Principles}

The mechanism must satisfy three constraints:
\begin{enumerate}[leftmargin=*]
    \item \textbf{Generality.} The mechanism must not rely on task-specific
    schemas or hand-crafted extraction rules. It should work across all LIDPs
    in the benchmark.
    \item \textbf{Timeliness.} Information must be captured \textit{when it is
    visible}, not retrieved retroactively after it is lost.
    \item \textbf{Goal-grounded.} What to remember should be determined by the
    agent's own reasoning about which facts will be needed later to achieve
    the goal.
\end{enumerate}

\subsubsection{Mechanism Overview}

We introduce the \textbf{Long-chain Implicit Decision Mechanism} (LIDM),
an explicit working memory system that operates in three phases at every
agent step:

\begin{enumerate}[leftmargin=*,noitemsep,topsep=2pt]
    \item \textbf{Goal-Aware Memory Prompt (GAMP).} Injected into the action
    selection prompt, instructing the LLM to identify and record facts from the
    current screen that will be needed later.
    \item \textbf{Working Memory Context (WMC).} Injects all previously recorded
    facts into the prompt so the agent can use them for downstream decisions.
    \item \textbf{Fact Parsing and Storage (FPS).} Parses the LLM's structured
    memory output and persists it into a \texttt{WorkingMemory} store.
\end{enumerate}

\paragraph{Goal-Aware Memory Prompt.}
At every step $t$, the action prompt $P_t^{\text{action}}$ is extended with:

\begin{lstlisting}[basicstyle=\small\ttfamily, breaklines=true,
                   frame=single, xleftmargin=10pt]
--- Goal-Aware Memory ---
Your overall goal is: {goal}
Look at the current screen carefully. Based on your goal,
is there any specific information here (numbers, names,
dates, file contents, rules, or decisions) that you will
need LATER to complete the goal?
If yes, record it using: [REMEMBER: key=value, key2=value2]
If nothing needs to be remembered, you can skip this.
\end{lstlisting}

The LLM answers this inline, before producing its action. Crucially, this
happens \textit{while the information is still visible on screen}, not
after the agent navigates away.

\paragraph{Working Memory Context.}
The \texttt{WorkingMemory} store $\mathcal{M}_t$ at step $t$ contains all
facts recorded in previous steps:

\[
\mathcal{M}_t = \bigcup_{i=1}^{t-1} \mathcal{F}_i,
\qquad
\mathcal{F}_i = \{(k, v) \mid \text{extracted from step } i\}.
\]

The formatted memory text is injected into the action prompt as:

\begin{lstlisting}[basicstyle=\small\ttfamily, breaklines=true,
                   frame=single, xleftmargin=10pt]
--- Your Working Memory ---
  [Step 3] symptom_count = 2
  [Step 5] booking_rule = >3 -> Specialist
Use these facts when making decisions.
\end{lstlisting}

\paragraph{Fact Parsing and Storage.}
After the LLM produces its output $O_t$, a regular expression extracts all
\texttt{[REMEMBER: ...]} blocks. Each comma-separated \texttt{key=value} pair
is parsed into a dictionary and appended to $\mathcal{M}$:

\[
\mathcal{F}_t = \texttt{parse\_re}(O_t, \texttt{\[REMEMBER:\s*(.*?)\]}).
\]

\subsubsection{Integration with PegaAgent}

LIDM is fully integrated into PegaAgent's \texttt{step()} loop. The prompt
assembly at each step is:

\[
P_t = P_t^{\text{T3A}} \oplus P_t^{\text{loop}} \oplus P_t^{\text{WMC}}
\oplus P_t^{\text{GAMP}},
\]

where $\oplus$ denotes string concatenation, $P_t^{\text{T3A}}$ is the base
T3A action prompt, $P_t^{\text{loop}}$ is the loop-detection notification
(if triggered), $P_t^{\text{WMC}}$ is the working memory context, and
$P_t^{\text{GAMP}}$ is the goal-aware memory instruction.

The \texttt{WorkingMemory} store is initialized at agent construction and
cleared on reset, ensuring no cross-task contamination. The mechanism adds
one lightweight LLM inference step per turn (the GAMP is part of the same
prompt, not a separate call), incurring negligible overhead.

\subsubsection{Why This Works for All LIDPs}

LIDM does not presuppose \textit{what} to remember, \textit{when} to remember
it, or \textit{when} to use it. The LLM determines all three dynamically:
\begin{itemize}[leftmargin=*,noitemsep,topsep=2pt]
    \item \textit{What:} The LLM identifies relevant facts based on the goal.
    \item \textit{When to record:} Every step triggers the GAMP, so any
    visible information can be captured.
    \item \textit{When to use:} The WMC is injected into every step's prompt,
    so recorded facts are always available for decision-making.
\end{itemize}

This design converts implicit, lossy memory (history summaries) into explicit,
structured facts that persist across app boundaries and step transitions---the
core capability required for long-chain implicit decision problems.
